\section{临床问题：从医学知识到 MultiMedQA}

有了序列视角，课程首先进入最接近语言模型原始形态的临床文本。问题并不是“模型能否背出医学事实”，而是能否在开放问题中组织科学共识、覆盖必要信息、避免危险遗漏，并把答案写成医生与普通用户都能使用的形式。

\subsection{为什么临床知识评估困难}

临床章节的分隔页之后，讲者引出 \emph{Large Language Models Encode Clinical Knowledge}。动机页展示医学知识的规模和复杂性：疾病、症状、检查、治疗与个体背景相互依赖，而且专业知识持续更新。单个 multiple-choice 分数可以测局部选择能力，却无法告诉我们一段长答案是否准确、是否遗漏关键危险信号、是否对患者真正有帮助。

\lecturefigure{slide-10-section-clinical.jpg}{章节切换：临床应用}{官方录像恢复幻灯片 V010；Stanford CS25 Lecture 17}

\lecturefigure{slide-11-clinical-knowledge-paper.jpg}{指定工作：Large Language Models Encode Clinical Knowledge}{官方录像恢复幻灯片 V011；Stanford CS25 Lecture 17}

\lecturefigure{slide-12-clinical-motivation.jpg}{临床动机：医学知识规模大、关系复杂且高风险}{官方录像恢复幻灯片 V012；Stanford CS25 Lecture 17}

\begin{knowledgebox}{背景概念：医学问答的三层难度}
第一层是事实召回，例如某药物的作用；第二层是把多条事实组合到病例上下文中；第三层是沟通与风险管理，包括识别缺失信息、表达不确定性和建议就医。选择题主要覆盖前两层的一部分，长答案才暴露第三层，但长答案也更难自动评分。
\end{knowledgebox}

\subsection{从专用模型到大型通用模型}

时间线页回顾 biomedical language model 的演进：早期方法通常在 PubMed 或临床文本上训练专用模型，随后大规模通用语言模型开始以 few-shot 或 instruction-following 的方式进入医学 benchmark。下一页用“What's missing?”提醒读者：规模和流畅性尚未回答医学可信度、长答案评价与使用者差异。

\lecturefigure{slide-13-biomedical-llm-timeline.jpg}{生物医学语言模型的发展时间线}{官方录像恢复幻灯片 V013；Stanford CS25 Lecture 17}

\lecturefigure{slide-14-clinical-whats-missing.jpg}{已有医学语言模型还缺少什么？}{官方录像恢复幻灯片 V014；Stanford CS25 Lecture 17}

读图时应先比较训练范式，而不是只比较参数量：domain-specific pretraining 用专业语料改变表示；few-shot prompting 主要改变上下文；instruction tuning 用任务集合训练模型遵循指令；Med-PaLM 的 instruction prompt tuning 则进一步学习医学任务所需的软提示。它们分别干预语料、参数或提示，不能统称为“医学微调”。

\begin{warningbox}{大模型会答题，不等于会行医}
benchmark 题目通常已经给出封闭问题和候选答案，真实临床却需要主动收集信息、处理冲突证据并承担后果。课堂只讨论知识问答与文本辅助的研究证据，没有证明模型能够独立诊断、开药或代替医生。
\end{warningbox}

\subsection{MultiMedQA：把分散基准放进同一框架}

Medical question answering 页先把任务分成专业考试、医学研究问答与消费者健康问题。MultiMedQA 随后汇集 MedQA、MedMCQA、PubMedQA、MMLU clinical topics 等多项数据，并加入 HealthSearchQA 一类真实消费者问题。这样做的价值不是制造一个万能总分，而是让模型同时面对封闭选择与开放长答案。

\lecturefigure{slide-15-medical-question-answering.jpg}{医学问答：从考试题到真实用户问题}{官方录像恢复幻灯片 V015；Stanford CS25 Lecture 17}

\lecturefigure{slide-16-multimedqa-benchmark.jpg}{MultiMedQA：多数据集医学问答基准}{官方录像恢复幻灯片 V016；Stanford CS25 Lecture 17}

\lecturefigure{slide-17-multimedqa-examples.jpg}{MultiMedQA 中不同来源的问题示例}{官方录像恢复幻灯片 V017；Stanford CS25 Lecture 17}

\lecturefigure{slide-18-multimedqa-summary.jpg}{MultiMedQA 的任务范围与评价层次}{官方录像恢复幻灯片 V018；Stanford CS25 Lecture 17}

\begin{knowledgebox}{读图：不要把不同数据集直接平均成“医学能力”}
MedQA 近似美国执照考试风格，MedMCQA 覆盖多学科医学选择题，PubMedQA 更接近文献证据判断，MMLU clinical topics 是通用评测中的医学子集，消费者问题则没有天然候选项。它们测量的知识来源、答案形式与评分可靠性不同；总览应被当作 coverage map，而不是单一潜变量的精确测量。
\end{knowledgebox}

从数据设计角度看，MultiMedQA 同时制造了一个机会和一个张力。机会是同一个基础模型可以在统一 prompting 接口下接受多种问题，研究者因而能观察能力是否跨数据集迁移；张力是不同数据集的 label semantics 并不一致：PubMedQA 的“yes/no/maybe”依赖给定文献语境，MedQA 的最佳选项依赖考试规范，消费者问题往往需要澄清信息而不是立即作答。若训练时把所有答案都当作同质 token loss，模型可能学会形式而不是决策条件。更稳妥的系统会保留 dataset identity、answer type、evidence requirement 和 abstention policy，并分别报告每个子任务，而非用一个宏平均数掩盖失败集中在哪类用户问题上。

\teachervoice{讲者强调，医学是 knowledge intensive 的领域，但“模型知道答案”还不够。真正缺失的是让专家与普通用户都能评价的长答案框架，因为医生关注科学共识、遗漏与伤害，患者还关心是否读得懂、是否回应了自己的意图。}

\subsection{本章小结}

MultiMedQA 的贡献首先是扩大问题类型和使用者范围。它把医学选择题、研究问答与消费者长问题放到同一研究框架中，同时暴露一个关键事实：越接近真实沟通，越不能只依赖 exact match 或多项选择准确率。

